\documentclass[11pt]{article}
\usepackage[margin=1in]{geometry}
\usepackage[T1]{fontenc}
\usepackage[utf8]{inputenc}
\usepackage{amsmath,amssymb,mathtools}
\usepackage{booktabs}
\usepackage{array}
\usepackage{enumitem}
\usepackage{graphicx}
\usepackage{hyperref}
\usepackage{xcolor}
\usepackage{siunitx}

\title{Pulse-Waveform Design, Phase-Engineering, and\\
Corrected-Parameterization Multitone Optimisation\\
for Selective Qubit Rotation in Dispersive cQED}
\author{}
\date{\today}

\begin{document}
\maketitle

\section{Objective and Scope}

This document is a unified report for the
\texttt{sqr\_pulse\_waveform\_design} study. It has three parts:

\begin{enumerate}[label=\arabic*.]
    \item A reproduced summary of the original study, reconstructed from the
    saved data, figures, and validation artifacts after the prior manuscript
    source was overwritten.
    \item A phase-compilation follow-up: whether the remaining strict-logical
    error is mostly a correctable cavity-only Fock-dependent phase profile or a
    deeper waveform/control limitation.
    \item A corrected-parameterization multitone optimisation study
    (\texttt{run\_simultaneous\_multitone.py}): a reformulation of the multitone
    ansatz with explicit per-tone amplitude, phase, and frequency corrections
    $(d_\lambda, d_\alpha, \delta\omega)$, optimised for both the true SQR and
    the conditional-phase SQR (controlPhaseSQR) targets.
\end{enumerate}

The baseline study established the envelope hierarchy, the practical operating
window, the role of GRAPE as an upper bound, and the qualitative importance of
phase structure. Part~II turns that diagnosis into a direct compiled-control
test. Part~III introduces a principled correction parameterisation that unifies
the single-branch and all-branch multitone cases under one framework.

\section{Part I: Reproduced Original Study}

\subsection{System, targets, and waveform families}

The original study used a dispersive transmon--cavity model with
$\omega_q/2\pi = \SI{6.150}{\giga\hertz}$,
$\omega_c/2\pi = \SI{5.241}{\giga\hertz}$,
$\alpha/2\pi = \SI{-255}{\mega\hertz}$,
$\chi/2\pi = \SI{-2.84}{\mega\hertz}$, a qutrit transmon
($|g\rangle, |e\rangle, |f\rangle$), and a truncated cavity logical window
$n=0,\dots,3$. The main scan variable was the dimensionless selectivity
parameter $\chi T / 2\pi$.

Four baseline families were compared in the closed system:
\begin{itemize}
    \item single-tone Gaussian,
    \item square,
    \item cosine-squared (Hann),
    \item one-segment multitone with a Gaussian envelope.
\end{itemize}

The original questions were whether true SQR or conditional-phase SQR is the
right primitive, which smooth envelope performs best, where the practical
selectivity--decoherence optimum lies, and how far structured parametric pulses
sit below the GRAPE upper bound.

\subsection{Closed-system baseline conclusions}

Figure~\ref{fig:baseline-core} reproduces the central original result. The
conditional-phase target is dramatically easier than true SQR, and smooth
single-tone envelopes already reach essentially unit conditional-phase fidelity
in the practical regime.

\begin{figure}[t]
    \centering
    \includegraphics[width=\textwidth]{../figures/fidelity_vs_chi_t.pdf}
    \caption{Reproduced baseline closed-system result from the original study.
    Smooth single-tone envelopes rapidly approach unit conditional-phase SQR
    fidelity, while true-SQR fidelity remains limited by spectator phases.}
    \label{fig:baseline-core}
\end{figure}

The saved baseline data give the following concrete numbers:
\begin{itemize}
    \item Gaussian conditional-phase fidelity rises from $0.555$ at
    $\chi T/2\pi = 0.5$ to $0.998$ at $2.0$ and $0.99999$ at $3.0$.
    \item Cosine-squared performs best at intermediate duration, reaching
    $0.9963$ already at $\chi T/2\pi = 1.5$, $0.9981$ at $2.0$, and
    $0.999997$ at $3.0$.
    \item Leakage across the full baseline scan stays below
    $7.4\times 10^{-5}$, validating the effective two-level picture for the
    qubit manifold.
\end{itemize}

The reproduced message is the same as in the original report: conditional-phase
SQR, not strict true SQR, is the natural baseline primitive.

\subsection{GRAPE, decoherence, and practical operating point}

The original study also established three system-level facts that remain intact:

\begin{enumerate}[label=\arabic*.]
    \item \textbf{GRAPE is the proper upper bound.}
    Saved GRAPE results give conditional-phase fidelity
    $F_{\mathrm{GRAPE}} \approx 1$ for all scanned
    $\chi T/2\pi \in \{1,2,3,5,7,10\}$, while true-SQR GRAPE is already
    $0.8295$ at $\chi T/2\pi = 1$ and exceeds $0.987$ by $\chi T/2\pi = 3$.

    \item \textbf{The practical open-system optimum is short and smooth.}
    In the original Phase-5 data, the cosine-squared pulse achieved the best
    net conditional-phase fidelity, peaking at $0.9845$ around
    $\chi T/2\pi = 1.5$. Gaussian and multitone one-segment baselines peaked
    near $\chi T/2\pi = 2.0$ at about $0.9814$.

    \item \textbf{The original hardware-constrained GRAPE study did not see a DAC
    or amplitude-ceiling bottleneck.}
    The saved hardware-constrained and unconstrained GRAPE conditional-phase
    fidelities are both numerically $1.0$ across the scan.
\end{enumerate}

These three points matter because they show the system is controllable enough
for high-fidelity logical action; the main remaining problem is identifying the
right \emph{structured} primitive below the GRAPE ceiling.

\subsection{Generalized targets, optimized multitone, and all-branch behavior}

The original extension phases refined the interpretation of that gap:

\begin{itemize}
    \item The generalized-target study showed that echoed sequences can improve a
    block-phase-relaxed metric while hurting strict logical performance via
    same-block loss and leakage.
    \item The optimized single-branch multitone study improved the block-relaxed
    fidelity substantially only at very short gate times; for
    $\chi T/2\pi \ge 3$ it provided at most a marginal gain over the smooth
    single-tone baseline.
    \item The all-branch study produced the sharpest warning sign: DE$+$NM
    multitone pulses reached $F_{\mathrm{block}} \ge 0.987$ but only
    $F_{\mathrm{strict}} \le 0.27$, while GRAPE stayed at $F \ge 0.9999$.
\end{itemize}

Figure~\ref{fig:baseline-allbranch} reproduces that last point.

\begin{figure}[t]
    \centering
    \includegraphics[width=\textwidth]{../figures/allbranch_fidelity_comparison.pdf}
    \caption{Reproduced original all-branch result. Parametric multitone pulses
    can be very good under a block-phase-relaxed metric while remaining very far
    from the strict logical unitary.}
    \label{fig:baseline-allbranch}
\end{figure}

This all-branch gap motivated the follow-up: was the missing error mostly a
cavity-only diagonal phase structure, or something deeper?

\section{Part II: Phase-Compilation Follow-Up}

\subsection{Reframed control objective}

The new follow-up was organized around the compiled primitive
\[
U_{\mathrm{compiled}} \approx U_{\mathrm{diag}} U_{\mathrm{SQR}},
\qquad
U_{\mathrm{diag}} = \sum_n e^{i\phi_n} |n\rangle\langle n|.
\]
The key question was whether a cavity-only phase layer, combined with one global
qubit virtual-$Z$, can recover most of the strict-logical gap left by the
baseline waveforms.

This defines three nested gauges:
\begin{enumerate}[label=(\alph*)]
    \item raw strict logical fidelity,
    \item strict fidelity after one global qubit virtual-$Z$ and cavity-only
    phase compilation,
    \item branch-local-$Z$-relaxed fidelity, which serves as the upper bound if
    the remaining defect is mostly per-branch qubit phase rather than bosonic
    phase.
\end{enumerate}

\subsection{Follow-up cases}

Two complementary follow-up cases were run.

\paragraph{Single-target enlarged-window case.}
The standard single-target $R_X(\pi)$ SQR on branch $n_0 = 1$ was re-evaluated on
an enlarged logical window $n=0,\dots,7$ ($N=8$, $n_{\mathrm{cav}}=10$) for the
Gaussian and cosine-squared families over
$\chi T/2\pi \in \{1,2,3,5\}$.

\paragraph{Representative short-gate structured all-branch case.}
The previously defined Phase-B independent-tone all-branch multitone optimizer
was re-run at $\chi T/2\pi = 0.5$ and $1.0$ to test whether cavity-only
compilation can rescue a structured multitone family precisely where the
original report suggested phase engineering might matter most.

\subsection{Single-target SQR does induce a cavity-only phase profile, and it is almost perfectly linear}

Figure~\ref{fig:phase-profile} shows the extracted cavity-only phase profile for
the enlarged-window single-target problem. The phase is not only present; it is
extremely smooth. Across both families and all scanned durations, the worst
linear-fit RMS phase error is only $4.59\times 10^{-5}$ rad.

At $\chi T/2\pi = 3$, both baselines are described to numerical precision by
\[
\phi(n) \approx -0.0684\,n \quad \text{rad},
\]
after gauge-fixing $\phi(0)=0$.

\begin{figure}[t]
    \centering
    \includegraphics[width=\textwidth]{../figures/phase_compilation_phase_profiles.pdf}
    \caption{Follow-up phase extraction on the enlarged $N=8$ logical window.
    The fitted cavity-only phase profile is essentially linear in photon number
    for both Gaussian and cosine-squared single-target SQR.}
    \label{fig:phase-profile}
\end{figure}

This is the strongest positive answer produced by the follow-up: in the
single-target setting, the cavity-only phase profile is real, low-complexity,
and naturally compilable.

\subsection{Cavity-only compilation helps on larger logical windows, but it does not close the full gap}

Figure~\ref{fig:single-recovery} summarizes the single-target compiled-control
comparison.

\begin{figure}[t]
    \centering
    \includegraphics[width=\textwidth]{../figures/phase_compilation_single_target_recovery.pdf}
    \caption{Follow-up compiled comparison on the enlarged $N=8$ logical
    window. A linear cavity-only phase layer improves the strict logical action,
    but does not reach the branch-local-$Z$-relaxed upper bound.}
    \label{fig:single-recovery}
\end{figure}

The concrete improvements are:
\begin{itemize}
    \item Gaussian: $0.738 \to 0.755$ at $\chi T/2\pi = 1$ and
    $0.927 \to 0.949$ at $\chi T/2\pi = 3$.
    \item Cosine-squared: $0.777 \to 0.795$ at $\chi T/2\pi = 1$ and
    $0.931 \to 0.954$ at $\chi T/2\pi = 3$.
\end{itemize}

The gain is small but systematic. The minimum improvement over the scan is
$0.0167$, and for the practical regime $\chi T/2\pi \ge 2$ it remains above
$0.0217$ for both families.

However, cavity-only compilation does \emph{not} recover the full gap to the
branch-local-$Z$-relaxed fidelity. At $\chi T/2\pi = 3$, the cosine-squared
pulse reaches $F_{\mathrm{diag}} = 0.954$ after cavity-only compilation, while
the branch-local-$Z$-relaxed upper bound is already $0.999998$.

So the single-target follow-up gives a nuanced answer:
\begin{quote}
The strict-logical deficit is partly a cavity-phase problem, but not purely a
cavity-phase problem.
\end{quote}

\subsection{Coherent superposition benchmarks improve after cavity-phase compilation}

The enlarged-window coherence benchmark confirms that the compiled phase layer is
not just improving an operator-overlap metric. At $\chi T/2\pi = 3$:
\begin{itemize}
    \item Gaussian pair-superposition fidelity improves from mean/min
    $0.959/0.671$ to $0.974/0.868$.
    \item Cosine-squared improves from mean/min $0.961/0.686$ to
    $0.976/0.879$.
\end{itemize}

\begin{figure}[t]
    \centering
    \includegraphics[width=0.92\textwidth]{../figures/phase_compilation_superposition_benchmarks.pdf}
    \caption{Follow-up coherent-superposition benchmark at $\chi T/2\pi = 3$ on
    the enlarged $N=8$ logical window. A cavity-only compiled phase layer
    improves both the mean and worst-case probe fidelity.}
    \label{fig:coherence}
\end{figure}

\subsection{The representative all-branch short-gate problem is not a cavity-only phase-compilation problem}

The same conclusion does not hold for the structured all-branch short-gate case.
The representative Phase-B independent-tone optimizer gives:

\begin{center}
\begin{tabular}{ccccc}
\toprule
$\chi T/2\pi$ & $F_{\mathrm{raw}}$ & $F_{\mathrm{diag}}$ & $F_{\mathrm{branchZ}}$ & $F_{\mathrm{GRAPE}}$ \\
\midrule
0.5 & 0.138 & 0.028 & 0.972 & 0.9999 \\
1.0 & 0.058 & 0.059 & 0.565 & 1.0000 \\
\bottomrule
\end{tabular}
\end{center}

\begin{figure}[t]
    \centering
    \includegraphics[width=0.9\textwidth]{../figures/phase_compilation_allbranch_short_gate.pdf}
    \caption{Follow-up all-branch short-gate comparison. Cavity-only bosonic
    phase compilation does not rescue the strict logical action, while
    branch-local-$Z$ freedom and GRAPE reveal a much larger reachable control
    set.}
    \label{fig:allbranch-short}
\end{figure}

At $\chi T/2\pi = 0.5$, the fitted cavity-only phase profile is irregular,
\[
\phi \approx (0,\ 1.50,\ -0.14,\ 1.36)\ \mathrm{rad},
\]
and applying it actually \emph{worsens} the strict logical fidelity. This means
the structured all-branch short-gate failure mode is not a missing bosonic
SNAP-like layer; it is a deeper branch-local qubit-phase and/or non-diagonal
control problem.

\section{Unified Interpretation}

The combined baseline-plus-follow-up picture is now much clearer.

\paragraph{What the original study already proved.}
Smooth single-tone selective pulses, especially cosine-squared, are already
excellent conditional-phase SQR primitives in the practical operating regime.
GRAPE demonstrates that the system itself is not the bottleneck. The all-branch
extension showed that phase structure matters and that block-relaxed success can
severely overestimate strict logical quality.

\paragraph{What the follow-up adds.}
The missing strict-logical error is not of a single type.
\begin{itemize}
    \item For single-target SQR on larger logical windows, there is a real and
    almost perfectly linear cavity-only phase profile. Compiling it helps.
    \item But even there, the residual gap to the branch-local-$Z$ limit is
    substantial, so cavity-only compilation is not the whole story.
    \item For structured all-branch short-gate multitone control, cavity-only
    compilation is the wrong correction class entirely.
\end{itemize}

So the strongest version of the hypothesis
``the remaining strict-logical error is mostly a bosonic diagonal phase''
is not supported uniformly across the study space. The correct statement is
narrower and more useful:
\begin{quote}
Low-order cavity-phase compilation is a real, calibratable improvement for
single-target SQR on larger logical windows, but the hardest structured
multitone failures are not bosonic-phase problems.
\end{quote}

\section{Practical Recommendation}

The unified experiment-facing recommendation is summarized in
Table~\ref{tab:decision}.

\begin{table}[t]
\centering
\small
\resizebox{\textwidth}{!}{%
\begin{tabular}{>{\raggedright\arraybackslash}p{3.5cm}ccccc>{\raggedright\arraybackslash}p{2.8cm}}
\toprule
Primitive & Nominal fidelity & Strict logical fidelity & Robustness & Hardware realism & Calibration complexity & Recommendation \\
\midrule
Single-tone cosine-squared SQR & $\approx 1.000$ conditional-phase / branch-local-$Z$ relaxed near $\chi T/2\pi = 2$--$3$ & High on the original $N=4$ window; $0.931$ on enlarged $N=8$ window at $\chi T/2\pi = 3$ & High in the original study & High & Low & Best immediate single-target primitive \\
Cosine-squared SQR + linear cavity compilation & Same pulse plus fitted $\phi(n)\approx -0.068 n$ on enlarged window & $0.954$ on enlarged $N=8$ window at $\chi T/2\pi = 3$ & High & High & Medium & Best immediate compiled upgrade when inter-Fock coherence matters \\
Optimized single-branch multitone & Helpful mainly at very short gate times; marginal gain for $\chi T/2\pi \ge 3$ & Similar to smooth single-tone in the practical regime & Medium & Medium & Medium & Secondary option; not a practical-regime winner \\
Structured all-branch multitone + cavity compilation & Can look good under branch-local/block gauges & Fails badly in representative short-gate cases ($0.028$ at $\chi T/2\pi = 0.5$) & Unknown & Medium & Medium-high & Do not deploy as a bosonic-compiled primitive \\
GRAPE pulse & $\approx 1.0$ & $\approx 1.0$ & Not yet fully mapped here & Medium & High & Long-term benchmark / control target \\
\bottomrule
\end{tabular}
}
\caption{Unified recommendation table for the reproduced original study and the
new phase-compilation follow-up.}
\label{tab:decision}
\end{table}

The practical decision is therefore:
\begin{enumerate}[label=\arabic*.]
    \item Use cosine-squared single-target SQR as the default experiment-facing
    primitive in the practical $\chi T/2\pi \sim 2$--$3$ regime.
    \item When logical windows are enlarged and inter-Fock coherence matters, add
    a low-order cavity-only phase compilation layer; it is simple, low-complexity,
    and demonstrably helpful.
    \item Do not interpret poor short-gate structured all-branch performance as
    evidence that ``only a SNAP layer is missing.'' In that regime the control
    limitation is deeper.
\end{enumerate}

% ===========================================================================
\section{Part III: Corrected-Parameterisation Multitone Optimisation}
% ===========================================================================

\subsection{Implemented Waveform Forms}

\subsubsection{Single-Tone Rotation}

A general single-tone qubit rotation has the form
\begin{equation}
w[k]
=
s\,w_0[k]\,
e^{i\phi_{\mathrm{eff}}}\,
e^{i\omega t_k},
\end{equation}
where
\begin{align}
\phi_{\mathrm{eff}} &= \phi + d_\alpha, \\
s &= \frac{\theta}{\pi} + \frac{d_\lambda}{\lambda_0(T)},
\qquad
\lambda_0(T) = \frac{\pi}{2T}.
\end{align}

The physical IQ outputs are the real and imaginary parts:
\begin{equation}
I[k] = \Re\{w[k]\},
\qquad
Q[k] = \Im\{w[k]\}.
\end{equation}

The scale $s=1$ with $d_\lambda=0$ and $\theta=\pi$ recovers exactly the
standard $\pi$-pulse amplitude $\lambda_0(T)=\pi/(2T)$. The shared base
envelope $w_0[k]$ is the area-normalised Gaussian
($\int_0^T w_0(t)\,dt = 1$), so the actual Rabi amplitude per tone is
$s\,\lambda_0(T) = \theta/(2T) + d_\lambda$.

\subsubsection{Multi-Tone Conditional Rotation}

A conditional rotation is implemented as a superposition of tones,
\begin{equation}
w[k] = \sum_{n=0}^{n_{\max}-1} w_n[k],
\end{equation}
with each tone given by
\begin{equation}
w_n[k]
=
s_n\,w_0[k]\,
e^{i\phi_{\mathrm{eff},n}}\,
e^{i\omega_n t_k}.
\end{equation}

Here,
\begin{align}
s_n &= \frac{\theta_n}{\pi} + \frac{d_{\lambda,n}}{\lambda_0(T)}, \\
\phi_{\mathrm{eff},n} &= \phi_n + d_{\alpha,n}, \\
\omega_n &= \omega_{\mathrm{det},n} + \delta\omega_n.
\end{align}

A tone is included in the waveform only when $|s_n|$ is non-negligible
($|s_n| > 10^{-4}$), rather than solely when $\theta_n \ne 0$.  This
allows the optimiser to ``activate'' spectator-branch tones via the
correction $d_{\lambda,n}$ alone, without needing to set a non-zero
nominal rotation angle for that branch.

\subsection{Free parameters and initial conditions}

Each tone $n$ contributes three scalar correction parameters:
\[
\text{free} = \{d_{\lambda,n},\; d_{\alpha,n},\; \delta\omega_n\}_{n=0}^{n_{\max}-1}.
\]
For $n_{\max}=4$ tones this gives $3\times 4 = 12$ free parameters.

The two study cases set the \emph{nominal} parameters as follows:

\paragraph{Single-branch SQR / controlPhaseSQR.}
Target branch $n_0=1$ has $\theta_{n_0}=\pi$, $\phi_{n_0}=0$; all
spectator branches have $\theta_n=0$, $\phi_n=0$.  At zero corrections
the pulse collapses to the standard single-tone Gaussian on branch $n_0$.

\paragraph{All-branch simultaneous rotation.}
All branches have $\theta_n=\pi$, $\phi_n=0$.  At zero corrections the
pulse is the common-envelope equal-amplitude multitone previously studied.

\subsection{Optimisation objectives}

Two objectives are optimised independently:
\begin{itemize}
    \item \textbf{SQR}: maximise the strict logical fidelity
    $F_{\mathrm{strict}} = |\mathrm{Tr}(U_{\mathrm{target}}^\dagger
    U_{\mathrm{restricted}})|^2 / d^2$ where $d = 2n_{\max}$.
    \item \textbf{controlPhaseSQR}: maximise the block-phase-relaxed
    fidelity $F_{\mathrm{block}}$, which allows an independent virtual-$Z$
    correction on each Fock branch.
\end{itemize}

Each optimisation uses a differential-evolution warm-start followed by
adaptive Nelder--Mead (12 parameters, up to 500 function evaluations per
Nelder--Mead run).

\subsection{Results and interpretation}

Results are stored in \texttt{data/simultaneous\_multitone\_results.npz}
after running \texttt{scripts/run\_simultaneous\_multitone.py}.  The output
contains the following key arrays for each of the four case--objective
combinations (prefixes \texttt{sb\_sqr}, \texttt{sb\_cphase},
\texttt{ab\_sqr}, \texttt{ab\_cphase}):

\begin{itemize}
    \item \texttt{*\_baseline\_*}: metrics with zero corrections (standard
    single-tone / common-envelope multitone reference).
    \item \texttt{*\_opt\_*}: metrics at the optimised corrections.
    \item \texttt{*\_opt\_d\_lambdas}, \texttt{*\_opt\_d\_alphas},
    \texttt{*\_opt\_delta\_omegas}: the optimised correction vectors
    (shape $n_{\chi T}\times n_{\max}$).
    \item \texttt{*\_opt\_s\_values}: the resulting scale vector $s_n$ at
    the optimum (indicates which tones were activated).
    \item \texttt{grape\_*\_fid}: GRAPE upper-bound fidelities at
    $\chi T/2\pi \in \{1,2,3,5\}$.
\end{itemize}

The key questions this study answers are:
\begin{enumerate}[label=\arabic*.]
    \item Does allowing explicit amplitude corrections $d_{\lambda,n}$ to
    activate spectator-branch tones improve the conditional-phase or strict
    SQR fidelity beyond the single-tone Gaussian baseline?
    \item For the all-branch target, do the three-correction-per-tone
    degrees of freedom close any of the gap previously identified between
    DE$+$NM multitone pulses and GRAPE?
    \item How do the optimised $s_n$ values distribute across tones, and
    which tones are actually activated by the optimiser?
\end{enumerate}

The $s_n$ vector provides an interpretable diagnostic: $s_n \approx 1$
for the target tone and $s_n \approx 0$ for spectators at the
single-tone baseline; non-zero spectator $s_n$ values reveal the degree
to which cross-talk tones are beneficial.

% ===========================================================================
\section{Open Problems}

The merged three-part report sharpens the next-step agenda:
\begin{itemize}
    \item Can the residual spectator branch-local qubit-$Z$ structure be
    reduced by echoed or two-layer selective pulses without sacrificing the
    favorable cavity-phase behavior identified in Part~II?
    \item Does the corrected-parameterisation multitone (Part~III) close the
    gap to GRAPE in the all-branch case, or does the structured-ansatz ceiling
    persist?
    \item How do the optimised correction vectors $(d_\lambda, d_\alpha,
    \delta\omega)$ correlate with the linearly-fitted cavity-only phase profile
    identified in Part~II?  Are they consistent with a single bosonic-phase
    layer, or do they require genuine per-tone amplitude modulation?
    \item How do the enlarged-window compiled results change under the more
    realistic hardware-transfer and expanded noise models?
\end{itemize}

\section{Conclusion}

The three study parts now fit together in a coherent hierarchy.  The baseline
(Part~I) established that smooth conditional-phase SQR is already strong in the
practical regime and that GRAPE provides a much higher ceiling than current
structured ans\"atze.  The phase-compilation follow-up (Part~II) showed that a
low-order cavity-only phase layer is genuinely useful for single-target SQR on
larger logical windows, but that the hardest structured short-gate failures are
not bosonic-phase problems.

Part~III introduces a cleaner parameterisation language for that gap. The three
per-tone corrections $(d_\lambda, d_\alpha, \delta\omega)$ expose, in a
single 12-parameter family, both the standard single-tone limit (all spectator
corrections zero) and the full common-envelope multitone limit (all $\theta_n$
equal), with a smooth interpolation in between.  The resulting optimised pulses
for both SQR and controlPhaseSQR targets provide the next quantitative reference
point against which compiled cavity-phase strategies and GRAPE can be compared.

The best experiment-facing primitive from Parts I--II remains
\emph{cosine-squared single-target SQR, optionally compiled with a low-order
linear cavity-only phase layer when coherence across a larger logical window
matters}.  Part~III targets the unresolved bottleneck: whether principled
per-tone amplitude and frequency corrections can further close that gap without
the full complexity of GRAPE.

\end{document}
